% Generated draft: requires mathematical and source review.
\documentclass[11pt]{article}
\usepackage[margin=25mm]{geometry}
\usepackage{fontspec}
\setmainfont{Latin Modern Roman}
\usepackage{amsmath,amssymb,mathtools,unicode-math}
\setmathfont{Latin Modern Math}
\usepackage{enumitem,needspace,xurl,hyperref}
\setlength{\parindent}{0pt}
\setlength{\parskip}{5pt}
\setlength{\emergencystretch}{3em}
\newcommand{\N}{\mathbb{N}}
\newcommand{\Z}{\mathbb{Z}}
\newcommand{\Q}{\mathbb{Q}}
\newcommand{\R}{\mathbb{R}}
\newcommand{\C}{\mathbb{C}}
\newcommand{\F}{\mathbb{F}}
\newcommand{\A}{\mathbb{A}}
\newcommand{\PP}{\mathbb{P}}
\newcommand{\E}{\mathbb{E}}
\newcommand{\eps}{\varepsilon}
\newcommand{\dd}{\,\mathrm{d}}
\begin{document}
\section*{UnsolvedMath ID: 30002456}
\textbf{Source problem code:} catalogue entry
\subsection{Definitions and mathematical statement}
\textbf{Research question.} Why are plant leaves generally flat? The supplied expanded formulation asks which growth-coordination and mechanical-feedback mechanisms explain why most plant leaves remain approximately flat despite heterogeneous cellular growth.

Here heterogeneous cellular growth means that growth can vary among cells. Growth coordination concerns relationships among these local growth processes, and mechanical feedback concerns interactions through which mechanical conditions influence growth. These descriptions identify the intended subject; the supplied record provides no mathematical definitions or governing equations for these mechanisms.

The requested target is an explanation of approximate leaf flatness under heterogeneous growth, incorporating growth coordination and mechanical feedback. However, the record specifies neither a class of leaves or growth models, a reference geometry, a measure and tolerance for approximate flatness, a developmental time interval, nor mechanical constitutive laws, initial conditions, or boundary conditions. It also supplies no population or measure interpreting ``most.'' Consequently, this is an explanatory research question rather than a fully specified quantified mathematical conjecture. A precise theorem, stability claim, or mechanism-selection criterion cannot be extracted without additional hypotheses; none is imposed here.
\subsection{Short English statement}
Explain how coordination of heterogeneous cellular growth and mechanical feedback could maintain approximate flatness in most plant leaves. The supplied question leaves the mathematical model, meaning of flatness, and scope of ``most'' unspecified, so this draft preserves an unresolved research target.
\subsection{Sources}
\begin{itemize}
\item Supplied local UnsolvedMath registry record: permanent ID \(30002456\), source problem code \texttt{catalogue entry}, title ``Mechanisms Maintaining Flat Plant Leaves,'' in the Oberwolfach Reports Open Problems collection. No individual proposer is supplied; the registry lists proposed year 2013.
\item Supplied report citation: \emph{Material Theories (2014)}, Oberwolfach Reports, DOI 10.4171/owr/2013/59, pp.\ 3433--3485. Original OWR record: \url{https://doi.org/10.4171/owr/2013/59}. The background describes these reports as records of mathematical workshops and their open-problem discussions. No problem-specific page or speaker attribution is supplied.
\item Supplied literature source labeled ``Primary leaf growth mechanics study'': \url{https://ems.press/content/serial-article-files/46487}. No author, article title, or more detailed bibliographic information is supplied for this link.
\item Literature assessment checked on \texttt{2026-08-21}: supplied dated curation, not independently verified. Its status label is ``open.'' Its assessment is that the extracted question is broad, that primary mechanics work identifies growth-coordination constraints and model behavior, and that this work does not provide a general biological explanation of leaf flatness. The supplied background requests independent verification of time-sensitive status and cited claims.
\item The registry labels the statement \texttt{corrected\_verified} and reports checking it against the official Oberwolfach report. That verification claim is supplied metadata, not an independent verification performed for this draft.
\end{itemize}
\end{document}
